\documentclass[12pt]{article}
\usepackage[utf8]{inputenc}
\usepackage[T1]{fontenc}
\usepackage{hyperref}
\title{Putoh Language Resources for K-12 Curriculum Development}
\author{Compiled by Claude Sonnet 5.5 (claude-sonnet-5-5) for Andr\'as Kornai}
\date{}
\begin{document}
\maketitle

\section{Introduction}
Putoh (ISO 639-3 \texttt{put}) is an endangered Austronesian language of Indonesia. Global Recordings Network places it in East Kalimantan with a population of about 11,000 and lists Abai and Pa Kembaloh among its dialects. Sources disagree: Joshua Project's people profile places Putoh speakers in Papua, which appears inconsistent with the other sources. Wikipedia notes the code is linked with Lun Bawang and that Putoh may be a dialect of Lundayeh. Very little material exists, and none of it is educational.

\section{Primary Educational Resources}
No primers, textbooks or school programmes were found.

\section{Linguistic and Reference Materials}
\begin{itemize}
\item Wikipedia, Lun Bawang language (mentions the Putoh code): \url{https://en.wikipedia.org/wiki/Lun_Bawang_language}
\item Wikipedia, Apo Duat languages (Putoh as possible member or Lundayeh dialect): \url{https://en.wikipedia.org/wiki/Apo_Duat_languages}
\end{itemize}
No dictionary or grammar was found.

\section{Religious and Community Resources}
\begin{itemize}
\item Joshua Project language page for Putoh: \url{https://joshuaproject.net/languages/put}
\item Joshua Project people profile (location information conflicts with other sources): \url{https://joshuaproject.net/people_groups/14488}
\end{itemize}

\section{Digital and Audio Resources}
\begin{itemize}
\item Global Recordings Network, Putoh language page: \url{https://globalrecordings.net/en/language/put}
\end{itemize}

\section{Curriculum Development Context}
There is no evidence of Putoh in schools or of a written standard. Any curriculum work would first require documentation and clarification of the relationship to Lun Bawang/Lundayeh and to Abai.

\end{document}